% STATUS: draft
% LAST_MODIFIED: 2026-07-26
\documentclass{internal-report}
\usepackage[UTF8]{ctex}

% 内部报告路径设置在 preamble 中
% 编译时需要将 internal-report.cls 放在 latex/templates/ 下
% 或使用 TEXINPUTS 环境变量

\title{2024 CUMCM C题「农作物的种植策略」—— 初始分析}
\author{MathModel Team}
\date{\today}

\begin{document}

\maketitle

\section{问题重述}

某乡村地处华北山区，拥有 1201 亩露天耕地（34 个地块，分平旱地、梯田、山坡地、水浇地 4 种类型）、16 个普通大棚、4 个智慧大棚。需要为该乡村制定 2024--2030 年（共 7 年）的农作物最优种植方案，在耕地类型约束、轮作约束（不重茬 + 三年豆类循环）、预期销量限制下，最大化种植总收益。

题目的核心挑战在于：
\begin{itemize}
    \item \textbf{三维决策空间}：年份 $\times$ 地块 $\times$ 作物 $\times$ 季次，变量规模大
    \item \textbf{多类约束耦合}：线性（面积上限） + 组合（轮作） + 软约束（分散度、最小面积）
    \item \textbf{不确定性递进}：Q1 确定性 → Q2 参数随机 → Q3 相关性 + 模拟
    \item \textbf{跨年依赖}：豆类三年一种、不重茬等约束导致方案在时间维度上耦合
\end{itemize}

\section{子问题拆解}

\subsection{问题 1：静态确定性优化}

假设所有参数（销量、成本、产量、价格）均保持 2023 年水平不变，求 2024--2030 年最优种植方案。两种超产处理方式：
\begin{enumerate}
    \item 超产部分滞销浪费（完全损失）
    \item 超产部分按原价 50\% 降价出售
\end{enumerate}

\textbf{建模建议}：混合整数线性规划（MILP）。决策变量为 $x_{t,i,j,s}$（第 $t$ 年第 $i$ 地块第 $j$ 作物第 $s$ 季的种植面积）。目标函数为总利润（销售收入 $-$ 种植成本），约束包括：面积上限、产量-销量匹配（分段）、不重茬（$0$-$1$ 变量乘积）、豆类轮作、地块-作物匹配等。

\textbf{参考方案}：SP-003 中 038 使用 LINGO 精确求解，变量规模约 $7 \times 54 \times 41 \times 2 \approx 3.1 \times 10^4$，在 LINGO Global Solver 能力范围内。

\subsection{问题 2：不确定环境鲁棒优化}

在 Q1 基础上引入参数不确定性：
\begin{itemize}
    \item 小麦/玉米预期销量年增长 $5\%$--$10\%$
    \item 其他作物预期销量年变化 $\pm 5\%$
    \item 亩产量年变化 $\pm 10\%$（气候影响）
    \item 种植成本年增长 $\approx 5\%$
    \item 蔬菜价格年增长 $\approx 5\%$，食用菌价格年下降 $1\%$--$5\%$
\end{itemize}

\textbf{建模建议}：两种路线备选——
\begin{enumerate}
    \item \textbf{蒙特卡洛 + 益损值矩阵}（038 路线）：生成多组候选方案（各参数取最优/最劣/均值等组合），在 1000 组蒙特卡洛采样状态下评估各方案，用乐观/悲观/期望值三准则选最优方案。
    \item \textbf{风险因子分解法}（234 路线）：从气象文献中提取干旱、寒潮等具体冲击事件的减产率，作为确定性因子修正参数，做多场景对比，用利润波动指标 $\phi = (W' - W)/W$ 评估方案抗风险能力。
\end{enumerate}

\textbf{综合推荐}：结合两者——先用蒙特卡洛做基准评估，再用风险因子（干旱/寒潮）做领域增强，展现方法的互补性。

\subsection{问题 3：替代性/互补性 + 相关性}

在 Q2 基础上，考虑：
\begin{itemize}
    \item 作物间的可替代性（如不同蔬菜间）和互补性（如豆类肥田效应）
    \item 预期销量与价格、成本之间的相关性
    \item 通过模拟生成数据，与 Q2 结果对比
\end{itemize}

\textbf{建模建议}：
\begin{itemize}
    \item 替代性/互补性：定义作物间的替代/互补系数矩阵（连续）或替代品集合/互补品集合（离散，234 路线）
    \item 相关性：使用 Copula 或多元正态分布生成关联的随机参数
    \item 将修正后的需求/产量模型嵌入 Q2 优化框架
\end{itemize}

\section{相关知识}

\subsection{知识库关联卡片}
\begin{itemize}
    \item \textbf{SP-003}：2024 C 题三种建模哲学对比（精确 LP vs DEGA vs GA+风险因子），是本项目最直接的方法论参考
    \item \textbf{MS-018}：MILP 大规模种植排程精确优化（含不重茬、豆类轮作的线性化技巧）
    \item \textbf{AL-003}：遗传算法用于复杂约束全局优化（234 路线）
    \item \textbf{AL-011}：蒙特卡洛 + 益损值矩阵的稳健性评估（038 路线）
    \item \textbf{AL-012}：DEGA 混合算法（DE + GA，063 路线）
    \item \textbf{AL-013}：领域风险因子分解法（文献驱动，234 路线）
    \item \textbf{PR-001}：约束处理优先级：编码法 $>$ 罚函数 $>$ 修复法 $>$ 拒绝法
    \item \textbf{PR-015}：模型建立三段式与编号步骤法
\end{itemize}

\subsection{参考论文}
\begin{itemize}
    \item \textbf{2024C038}：精确 LP 路线（LINGO），Q2 蒙特卡洛益损值矩阵
    \item \textbf{2024C063}：DEGA 混合路线（DE + GA），三次样条适应度拟合
    \item \textbf{2024C234}：GA + 风险因子路线，替代品/互补品离散集合
\end{itemize}

\subsection{知识空白}
\begin{itemize}
    \item GAP-008：混合整数规划在 MCM 问题三中的通用建模模式
    \item 跨年动态规划 vs 逐年独立优化的对比分析（SP-003 指出的共同盲区）
\end{itemize}

\section{初始建模假设}

\begin{enumerate}
    \item \textbf{每年独立决策}：暂不考虑跨年策略博弈（如今年多种豆类以换取来年高产），每年优化独立进行。这是三篇论文的共同做法，也是 O-1 奖论文的隐含选择。
    
    \item \textbf{价格外生}：农作物的销售价格为外生给定，不受该乡村产量影响（小农假设，对市场价格没有定价权）。
    
    \item \textbf{确定性等价}：Q1 中所有参数用 2023 年的实际值作为未来 7 年的固定值；Q2/Q3 中参数按题目给定的变化率更新。
    
    \item \textbf{种植面积连续}：将种植面积视为连续变量（可分割到小数），避免引入过多 $0$-$1$ 变量。实际操作中可四舍五入。
    
    \item \textbf{季次定义}：第一季为春夏作物（3--8月），第二季为秋冬作物（9--次年2月）。具体季次划分需参考附件 1。
    
    \item \textbf{分散度和最小面积}：作为软约束处理，在目标函数中添加惩罚项，或在求解后做方案筛选。具体阈值通过灵敏度分析确定。
\end{enumerate}

\section{待解决的开放问题}

\begin{enumerate}
    \item \textbf{附件数据读取}：需要从附件 1 和附件 2 的 Excel 中提取地块类型、作物-地块匹配矩阵、2023 年种植数据、作物收益率等关键参数。
    
    \item \textbf{分散度/最小面积的量化}：题目只定性要求"不能太分散""面积不宜太小"，需要定义可操作的量化阈值。038 使用面积占比 $\geq$ 某一比例（如 25\%）和地块数上限，但该阈值本身需要灵敏度分析。
    
    \item \textbf{求解器选择}：LINGO（精确 LP） vs Python + Gurobi/CVXPY（灵活） vs 遗传算法（复杂约束）。需要权衡精度、开发效率和模型灵活性。
    
    \item \textbf{Q2 不确定性建模路线}：蒙特卡洛益损值矩阵 vs 风险因子分解法，还是两者结合？需要根据题目对"种植风险"的定义做选择。
    
    \item \textbf{Q3 的模拟数据生成}：如何在满足相关性约束的前提下生成合理的模拟数据？Copula 方法是否过于复杂？
    
    \item \textbf{跨年方案的一致性}：虽然每年独立优化，但豆类三年一种约束使得方案在时间上耦合。如何处理"时序可行性"？
    
    \item \textbf{大鹏内部作物搭配}：普通大棚"一季蔬菜 + 一季食用菌"和智慧大棚"两季蔬菜"的具体搭配规则需要在附件 1 中确认。
\end{enumerate}

\section{建议的下一步}

\begin{enumerate}
    \item \textbf{数据探索}：用 Python 读取附件 1 和附件 2，提取关键参数，生成地块-作物可行性矩阵。
    \item \textbf{先攻克 Q1}：Q1 是确定性优化，是 Q2/Q3 的基础。建议先用 MILP 建模 Q1，对比"浪费"和"半价"两种超产处理的利润差异。
    \item \textbf{方法论选择}：在 Q1 求解经验基础上，决定 Q2 采用哪种不确定性建模路线（蒙特卡洛 vs 风险因子 vs 两者结合）。
    \item \textbf{知识吸收}：精读至少一篇 2024C 参考论文（建议先读 2024C038，其方法最系统），提取可复用的建模技巧和参数设置。
\end{enumerate}

\end{document}
